\chapter{生态学经典论文}
\Large\textbf{Foundations of Ecology: Classic Papers with Commentaries}
\par \emph{Leslie A. Real, James H. Brown (Eds.)} \normalsize

\section{奠基性论文}

\subsubsection*{Sharon E. Kingsland. Defining Ecology as a Science}

生态学源于将实验方法、数量方法应用于有机体与环境关系，群落结构和演替，种群动力学的研究。背景是马尔萨斯人口论和达尔文的自然选择论\footnote{自然选择论受到了马尔萨斯的启发：达尔文推论过度生育会导致种内竞争加剧、死亡率上升，即选择过程。}。 达尔文的重要贡献是提出种内竞争是进化的主要动力\footnote{先前人们已认识到种间斗争的作用，如捕食者清除猎物中的弱势个体，但他们认为这只是通过清除异常个体保存了猎物的正常特征。而达尔文认识到种内竞争可以是创造性的而非保存性的。}。此外，Herbert Spencer认为自然界处于动态平衡中，启发生态学家探索自然具有自我调节能力、使种群间保持平衡的原因\footnote{Spencer接受自然选择观点，但认为生物会响应环境而变异，这一观点称为新拉马克主义。}。

1860年代Ernst Haeckel提出oecology的说法，用于描述有关生存斗争(struggle for existence)的研究。Ecology一词最初于1890年代被美国植物学家使用，用于摒弃描述性方法，关注有机体与环境关系的生理学研究\footnote{他们将生态学视为一种“户外的生理学”，继承生理学在实验室中的严格方法。受达尔文影响，关注是否可通过改变环境控制进化，这有显见的农业应用。}；动物学家紧跟这一趋势并于1920年代后来居上。

核心概念与方向：

\textbf{群落演替}：Henry Chandler Cowles通过假设空间变化能反映时间演替，在印第安纳沙丘地区首次得到了群落演替序列。他认为演替趋向稳定平衡态，但平衡态尚未达到，过程中也可能倒退。Frederic Clements认为演替的终点是顶级群落，而顶级群落的构成由优势植物确定。他还将植物群落类比为复杂有机体，还构建了植物群系的分类体系，二者后来都被弃用。在群落有机体隐喻的批评中，A. G. Tansley (1935)认为群落演替代表了有多个平衡态的动力系统的轨迹；并基于此定义了生物群落与物理环境构成的生态系统(ecosystem)；Henry Allen Gleason提出个体主义(individualistic)观点，认为每个植物群落是特定时空环境因素的独特产物，不存在固定不变的plant structure.

\textbf{食物网与生态位}：Stephen Forbes分析了种群中昆虫、鱼类、鸟类的食物关系，提出捕食、竞争关系调整种群增长率以维持群落平衡；其出生率与死亡率平衡的假设被长期沿用。Charles Elton定义了生态位的概念，用与描述物种在食物网中的位置，亦即群落中的“经济地位”。竞争排斥原理(competitive exclusion principle)指出不同物种不能占据相同生态位（为了共存，它们必须找到划分资源的方式）。这一原理此前已有观察(如共存物种回避竞争)。Elton并未论及这一原理, 而Joseph Grinnell将竞争排斥与生态位联系起来. Georgii F. Gause的工作进一步将这一原理推向中心地位: 种间竞争迫使两个物种生态位分离，这构成群落稳定性的基础。

\textbf{种群动力学}：Raymond Pearl提出了单一种群规模增长的logistic曲线，事实上Pierre-Francois Verhulst一个世纪前已有这一结果，只是未引起关注。Vito Volterra研究了种间竞争和捕食的数学模型，后者被称为Lotka-Volterra模型。Alfred James Lotka独立提出这一模型，其使用的联立微分方程类似后来Bertalanffy的一般系统论。Alexander John Nicholson和Victor Albert Bailey在种群模型中考虑了种内竞争、年龄结构等因素，不同于Lotka-Volterra，他们的模型得到震荡不断增大的不稳定系统\footnote{May指出对于小型、短生命物种（如昆虫与寄生虫），他们的模型比Lotka-Volterra模型更适用。}。

\textbf{能量流动}：Lotka研究了生物圈的能量转化，其著作Elements of Physical Biology涉及水、二氧化碳、氮、磷的循环。受Tansley的生态系统概念影响，Raymond Lindeman关注营养动力学(trophic-dynamic)视角，以湖泊为例研究生态系统中的能量流动及其对演替的影响。

\section{理论进展}

\subsubsection*{Leslie A. Real, Simon A. Levin. The Role of Theory in the Rise of Modern Ecology}

1950年代以来理论生态学主要围绕六个问题：

\textbf{群落的物种数量及丰度分布}：定量模型的构建始于1940年代, Fisher等 (1943)提出群落内个体数量为$k$的物种数为$\frac{\alpha \lambda^k}{k}$, 从而总物种数
\begin{equation}
    S = \sum_{k=1}^\infty \frac{\alpha \lambda^k}{k} = -\alpha \ln(1-\lambda)=\alpha \ln(1+\frac{N}{\alpha})
\end{equation}
其中$N=\frac{\alpha\lambda}{1-\lambda}$为个体总数. 但个体数量的众数为1, 与实际的单峰分布不符。

Preston的工作 (1948,1962)为这一方向基础: (1)对物种的个体数量取log bin得到正态分布曲线, 亦即个体数量服从对数正态分布；(2) 物种数量-面积满足的关系\footnote{这一关系基于物种相对丰度典范分布的性质。考虑每个log bin的个体总数得到的个体曲线, 其峰值往往对应物种曲线的尾部, 称满足这一性质的物种相对丰度分布为典范分布。}: $S=c A^z$, 实证给出$z=1/4$。

Preston没有给出对数正态分布的生成机制。这方面有两个角度：（1）生物角度：MacAuthur (1957,1960)提出“断棒(broken-stick)模型”，假设生态位空间表示为一维有限区间，同时选取$S-1$个断点，每段的长度对应一个物种的丰度。但这一模型并不能产生对数正态分布。Sugihara (1980)改进了这一模型，多维的超体代替一维区间，且断点逐一产生，得到了对数正态分布和物种数量-面积的$1/4$次幂关系。（2）统计角度：May (1975)认为对数正态分布是中心极限定理作用于多个因子交互的结果\footnote{May还提出，典范分布的性质也是大量物种对数正态曲线下的统计结论。}。

\textbf{生态位理论}：生态位的定义并不统一，归功于Grinnell, Elton, Gause和Hutchinson. Grinnell (1917,1928)侧重场所概念，即物种在环境中的位置，反映对物理环境的适应性；Elton (1927)侧重功能概念，即物种与食物、天敌的关系；Gause (1934)把Elton的观点扩展到种间竞争。Hutchinson (1957)基于集合论语言的形式化描述综合了场所、功能两种视角：设某一物种受$n$个环境变量$x_1,\dots,x_n$影响, 每个变量有其可接受区间；其笛卡尔积得到的超立方体代表了物种的基本生态位$N$(fundamental niche)。物种在现实中的场所$B$使$N$得到满足。不同物种基本生态位的交集反映了种间竞争。由于同一场所下不同物种的生态位不能相同，种间竞争影响下实际占据的生态位(realized niche)会不同于基本生态位。

\textbf{单一物种数量增长}：进化视角源于David Lack (1954)对鸟繁殖力的探讨，他指出鸟并未最大化产蛋数量\footnote{首先，一窝的产蛋数很低；其次，一个蛋被移走或破坏时，往往产一个蛋补上。这体现了生物生活史中的个体选择(individual selection)。}，而是最大化能存活的后代数。Cole (1954)进一步讨论了种群数量增长与生活史现象的联系，提出了一个悖论：单次生殖物种进化为多次生殖物种对种群数量的影响等价于每次繁殖后代数加1，而后者的进化难度远小于前者。后续研究给出了悖论的解决之法：繁殖与存活存在权衡，而幼体存活具有不确定性；只有成体、幼体存活率之比为1时悖论才成立。

动力学视角中里程碑研究来自May (1974). 在此之前，logistic增长模型表明种群数量会平滑地增加至最大承载量；密度过大对繁殖、存活的影响又会使其达到平稳值，这无法解释种群数量的实际波动。学界多尝试在连续增长模型基础上加入复杂因素，而May指出简单的离散增长模型即可产生复杂的动力学\footnote{Lorenz更早发现了非线性混沌系统，但当时影响有限；May的研究促使更多领域寻找时序中的混沌现象。}。假设代际不重叠，离散增长模型可表达为
\begin{equation}
    N_{t+1} = N_t \left(1+r\left(1-\frac{N_t}{k}\right)\right) \quad \text{或} \quad  N_{t+1} = N_t \exp\left(r\left(1-\frac{N_t}{k}\right)\right) 
\end{equation}
当$r$增大到某个大于2的临界值时，系统具有混沌性质：存在周期为任意正整数的循环，是否落入、落入何种循环完全依赖初值；不落入循环的初值有无穷多，且其序列性质类似随机过程。

\textbf{个体行为}：自然选择会偏好高效获取和利用资源的生物。MacArthur和Pianka (1966)使用最优化方法首次显式建模了资源获取策略，优化目标对应进化和行为假设。他们考虑的目标是最大化单位时间获取的期望能量；决策变量包括不同食物数量、觅食路径中包含的地块类型。Schoener (1971)提出需考虑另一种目标：最小化获取单位能量所需时间。Charnov (1976)针对地块选择问题，将每个地块停留时间作为决策变量。 

承接上述框架，捕食生态学后续研究主要关注生物个体行为的心理和经济学基础，包括引入营养约束；环境随机性与风险；学习、记忆、时间感知等心理约束；阐明能量获取与传统适应度（繁殖与生存）的联系；讨论行为的基因基础等。个体行为对种群增长、群落组织的影响，是一个关注较少的重要方向。

\textbf{多物种交互}：种间关系的数学建模始于1910年代Ronald Ross和W. Thompson对寄生虫病的研究；而成熟的框架归功于Lotka (1925)和Volterra (1926)，他们使用联立微分方程组表达不同物种数量动态。然而这一模型主要用于理论推演，未能成功标定并描述真实动态。

一些学者（如Margalef）认为群落物种越多稳定性越强，但仅停留于模糊论述，缺乏任何严格的根基。May (1975)使用Lotka的数学建模方法给出了完全相反的结论，动摇了前述观点，尽管其假设存在讨论余地。

\textbf{移动与尺度}：生物的散布(dispersal)体现了对环境时空异质性的响应。Skellam (1951)使用扩散模型建模生物的散布，尽管随机运动假设不实际，模型足以解释生物传播的宏观特性；考虑空间异质性可进一步提升预测能力。模型应用于外来物种入侵、赤潮\footnote{只有条件适宜的区域超过临界值(critical patch size)，微生物的增殖超过海水的扩散，赤潮才可能爆发。}等问题。

\section{争论与综合}

\subsubsection*{Joel G. Kingsolver and Robert T. Paine. Conversational Biology and Ecological Debate}

生态学经历着提出假设、提出不同假设、二者展开辩论的过程；有时单方面胜利，更常见的是问题重新定义后二者可以整合。作者认为已有讨论或未来将有的四个方向：(1)空模型与零假设的选取；(2)在通常关注的时空尺度，生态系统是否处于平衡态？一种观点是远离平衡态。许多环境存在干预，但其对生态系统稳定性的一般影响尚不清晰。(3)种间关系研究往往只关注二元关系，如何拓展以描述食物网的复杂结构？(4)种群和群落生态学往往只关注平均个体，个体差异仅仅是噪声吗，或者对种间交互有实际影响？

Hutchinson的论文\emph{Homage to Santa Rosalia} (1959)整合了Lack的进化生态学理念、生态位概念（种间竞争导致性状偏移），构建了生物多样性的理论。其基础有二：(1)生态位多样化产生物种多样性，提出性状参数1:3的差异可以使两物种共存；(2)复杂食物网使群落更稳定（后来May否定了这一论断）。

对于种群数量、群落结构的主要决定因素，有“种间竞争”和“非生物因素”两种观点。Hairston, Smith和Slobodkin (1960) 整合了两种观点，指出主要决定因素与考察的营养级有关，强调不同营养级的主导过程存在差异。

Ehrlich和Raven (1964) 由食草昆虫（蝶类）与植物的研究提出了协同进化(Coevolution)的概念：植物产生新的次级代谢物以抵御昆虫，食草昆虫进化出绕开防御的方式。特定昆虫与植物长期关联，表明协同进化能促进昆虫和植物的多样化。

Harper (1967) 提倡将动物进化生态学的理念引入植物学研究。

\section{方法进展}

\subsubsection*{James H. Brown. New Approaches and Methods in Ecology}

\textbf{古生态学}：Lennart von Post (1916) 指出泥炭沼泽的厌氧沉积物(anaerobic sediments)中可提取保存完好的花粉粒，结合沉积物的层次序列可反映最后一次大冰期以来（约11500年前至今）的植物物种变化，进而推断环境演变. 这一工作成为古生态学的基础——沼泽中的花粉粒、鼠类粪便堆中的植物组织、湖泊和海洋沉积物中的浮游生物骨架保存了近两万年的物种分布、群落组成信息。古生物学对古生态学影响有限，原因是化石往往是单一物种，重建远古群落和环境困难；其与当今生物的差异使得难以推断其生理、行为、生态关系。相对地，花粉粒保存了群落中大量物种，且与现存植物有较高的一致性。

\textbf{种群动力学}：Patrick Leslie (1945) 基于矩阵代数构建了生命表法(life tables)的数学基础：给定任意年龄$x$到$x+1$的存活率，及年龄$x$产生的平均后代数，可计算年龄的平稳分布；给定任意初始年龄分布，可计算人口增长率、到达平稳分布的时间\footnote{一些动物个体对种群数量的贡献依赖其社会阶层，为此可用分阶层的繁殖率、死亡率取代年龄的区分。}。先前Pearl、Lotka、Volterra的人口动力学模型面临参数难以标定的问题，而Leslie的方法依赖的数据是实验室和田野观测可获取的。Leslie的方法还应用于寿险精算。L. Charles Birch (1948)在实验室环境研究了象鼻虫的种群增长，使用生命表法估算了自然增长率等参数，发现温度等环境因素对自然增长率有显著影响。

C. S. Holling (1959) 在实验室和田野环境研究了小型哺乳动物对松叶蜂蛹的捕食行为。其最重要贡献是将捕食分解为两个层次：一是捕食者个体对猎物的搜索和摄食行为；二是影响捕食者和猎物种群数量的繁殖和死亡过程。这两部分对捕食量（作为猎物密度的函数）的影响分别称为功能响应(functional response)和数量响应(numerical response)。后续最优捕食策略的研究呼应了功能响应的理念。

\textbf{生理生态学}：Porter和Gates (1969) 研究生物体与环境能量交换的物理学，使用数学模型揭示太阳辐射、风速等环境因素对能量交换的影响；使用动物的金属模具测量物理因素的影响，体现生物过程（新陈代谢、皮肤蒸发等）的作用；提出气候空间(climate space)概念以描述生物体能维持热平衡从而生存、繁殖的物理因素（即基本生态位的物理维度）。后续研究持续关注物理环境因素对生物体生存、繁殖的影响。

\textbf{植物群落生态学}：Bray和Curtis (1957) 发展了一种基于物种组成对植物群落分类的定量方法，采用类似PCA的多变量分析方法。

\textbf{生态系统生态学}：1960年代，美国生态学界分为两派：Robert MacArthur领导的进化生态学，关注种群增长、群落组成、适应性等；Eugene Odum领导的生态系统生态学，关注能量流动、营养循环等。Odum (1969)的论文\emph{The strategy of ecosystem development}将两派关注过程统一于群落演替的概念框架中。进化生态学当时并不认可，因其似乎需要作用于群落或生态系统层次的自然选择过程；而后来有研究挑战了自然选择只作用于个体的正统观念，认为群落也存在一定的进化策略。论文同时为生态系统生态学提供了可检验的预测，推动这一分支走向实验科学。Odum还论及人类生态学，指出为了稳定性干预生态系统存在危险（如阻止大沼泽地区周年的洪水）。

1975年前尚可认为生态学是一门低技术学科，如今则不然，如计算机模拟；粒子计数器用于微生物计数；稳定同位素追踪化学物质的种间、空间流动；遥感用于土壤和植被分布探测；自动气相色谱(automated gas chromatography)识别植物的防御性化学物质；地理信息系统等。

\section{自然系统中的案例研究}

\subsubsection*{Robert K. Peet. Lessons from Nature}

物理学可以从极少的观测出发得出普遍结论，尽管观测的实现条件可能很特殊（如粒子对撞机）。一般而言，生物系统比物理系统更复杂（部分由于生物承载的遗传信息），生物学家选取有代表性的模式生物做透彻研究。生态系统更加复杂，除个体遗传信息外，还依赖个体空间分布和物理环境的特性。因此生态系统中的普遍性（generalization）依赖大量系统的比较研究。经典研究分为两类：一是单一系统的深入研究，产生新的洞察；二是在不同地点观测特定过程，得到广泛适用的结论。

\textbf{种群数量调控}：早期种群动力学模型基于种内竞争等密度相关（density dependent）因素，难以符合观测数据。Andrewartha、Birch和Davidson针对蝗虫和蓟马（1948）数量的系列研究\footnote{集成于专著\emph{The Distribution and Abundance of Animals}（1953）.}，揭示了种群数量调控中的密度无关因素。其中蓟马的研究基于14年的观测数据，发现30天前的环境变量可以解释78\%的种群数量对数波动。

\textbf{生态系统能量流}：Lindeman (1942)更偏向理论而非实证研究。Howard T. Odum（Eugene Odum之弟）针对Silver Spring, Florida的研究(1957)是自然系统能量流动的第一项相对完整的记录。Silver Spring的营养结构相对简单，边界清晰。结果中出人意料的是分解者的重要作用。John Teal先后在马萨诸塞州冷泉 (1957)和Sapelo Island, Georgia的盐沼 (1962) 开展了类似研究。在营养级间能量流动基础上，盐沼研究还量化了生物功能分组间的能量流动。意外结果包括：植物呼吸消耗了70\%以上能量；近一半的净初级生产被细菌消费；相当部分净初级生产因潮汐流出系统。

\textbf{植物群落结构}：为解决“群落有机体”的争论，Robert Whittaker开展了大烟山植被分布的研究，其结果支持Gleason的个体主义观点，宣告了Clements“群落有机体”观念的终结。Curtis和McIntosh更早发表了类似的案例研究（南威斯康星州植被），但因缺少对其结果implication的发掘和传播，影响力不及Whittaker。

Margaret Davis (1969)在Rogers Lake, Connecticut的孢粉学案例研究，改变了前人研究仅依据花粉占比推测沉积年代植被组成的方式，通过量化花粉汇集速率(influx rate)实现物种丰富度的更精确推断。基于花粉与现存物种差异的观察，Davis后来收集全北美花粉记录并通过制图反映不同植物在后冰期的不同传播速率和模式，为Gleason的个体主义观点提供了强力支撑。

\textbf{过程研究}：A.S. Watt的报告(1947)总结了自己的多个案例研究（覆盖沼泽、森林、苔原、草地多种类型），指出在成熟植物群落中仍存在小尺度的循环演替，更新了顶级群落同质、稳定的固有观念\footnote{此前也有类似观点，但仅针对特定类型的群落，Watt指出了这种再生过程的普遍性。}。事实上，植物群落具有时空动态性，地块可能处于不同发展状态。从倒木到森林大火的干扰事件都可视为动态的一部分，干扰事件的频率和规模分布反映不同生态系统动态，这成为当代生态学方向之一。Watt的观点还影响了后来发展的森林模拟模型，以及有关空间模式与异质性的研究。

MacArthur (1958)研究了云杉林中两种相似莺类得以共同存在的原因。其提倡构建理论-提出假设-检验假设的路径，推动了生态学的研究范式转变。

Brooks和Dodson (1965)观察到康涅狄格湖泊中外来鱼类使大型浮游动物消失的现象（无外来鱼类的湖泊中，大型浮游动物本来具有竞争优势）。这一案例研究表明捕食者可以控制其食用营养级的结构，催生了后来的营养级联(trophic cascade) 理论。
